% ============================================================
\chapter{Introduction}
\label{ch:introduction}
 
This chapter introduces the Automated Diagram Digitization project, presenting its motivation, problem context, objectives, and scope.
 
In today's digital era, the need to convert hand-drawn or printed diagrams into editable digital formats has grown across education, software development, and business process modelling. Flowcharts are frequently created on paper or whiteboards during brainstorming and design phases, yet they often remain in physical form, limiting shareability and editability. Traditional digitization involves manual redrawing in tools such as Microsoft Visio, Lucidchart, or draw.io—a time-consuming and error-prone process. Advances in computer vision and large language models have opened new possibilities for automated diagram recognition, enabling intelligent systems to understand and reconstruct diagrammatic information automatically.
 
%---------------------------------------------------------------
\section{Problem Statement}
%---------------------------------------------------------------
 
The transition from physical brainstorming artefacts to structured digital documentation remains a bottleneck in the software development lifecycle. Current workflows require developers and designers to spend 15–30 minutes recreating each diagram, introducing transcription errors and inconsistencies. The existing digitization ecosystem suffers from several interrelated challenges: manual redrawing bottlenecks, a total absence of automated shape and connection recognition, failure of standard OCR tools to preserve handwritten text or line breaks within nodes, inability to detect directional relationships, limited input flexibility, and restricted export options. These limitations collectively hinder the seamless transition from visual sketches to maintainable digital assets.
 
%---------------------------------------------------------------
\section{Objectives}
%---------------------------------------------------------------
 
The project aims to overcome traditional digitization limitations by delivering an AI-driven platform that covers the full flowchart lifecycle from recognition to export. Specific objectives include: developing an AI-powered recognition system using Google Gemini with custom prompt engineering for accurate shape, text, and connection detection; providing multiple input methods (image upload, text-to-flowchart conversion, and JSON import); implementing an interactive diagram editor with drag-and-drop, undo/redo, auto-layout, node and edge colour customisation, and bold/italic/font-size text formatting; supporting edge labelling via double-click; enabling export to PNG, SVG, and JSON; offering dark/light theme switching with preferences persisted via \texttt{localStorage}; and ensuring persistent flowchart storage using a PostgreSQL database with UUID-based identification and timestamp tracking.
 
%---------------------------------------------------------------
\section{Scope of the Project}
%---------------------------------------------------------------
 
\subsection{Functional Scope}
 
The system transforms raw visual inputs into interactive digital representations while preserving the original logical structure of diagrams. It identifies the four fundamental flowchart shapes from uploaded images, extracts text with multi-line line-break preservation using \texttt{\textbackslash n} characters, detects directional edges with appropriate handle assignments (top, bottom, left, right), and accepts images in PNG and JPG formats. Text-to-flowchart conversion transforms natural language algorithm descriptions into structured diagrams. JSON import and export enable portability. The interactive editor supports drag-and-drop from a shape palette, node resizing with adaptive text, undo/redo with keyboard shortcuts, auto-layout, edge labelling, and colour and text styling.
 
\subsection{Technical Scope}
 
The backend is implemented in FastAPI (Python) with asynchronous endpoints. The frontend uses React 19 with Vite and React Flow with custom node components. Integration with the Google Gemini API uses carefully engineered system prompts. A PostgreSQL database with JSONB support stores the complete React Flow state, using UUID primary keys and timestamp tracking. Export is implemented via \texttt{html-to-image} for PNG, React Flow's SVG renderer for vector output, and JSON serialisation for data portability.
 
\subsection{Limitations and Exclusions}
 
The project does not cover BPMN symbols, swimlanes, UML diagrams, or real-time video processing. It focuses on the four basic flowchart shapes commonly used in algorithm and process documentation. Support for additional diagram types, collaborative editing, and mobile interfaces are identified as future directions, discussed in Chapter~8.
 
%---------------------------------------------------------------
\section{Significance of Work}
\label{sec:significance}
%---------------------------------------------------------------
 
The platform addresses a genuine problem faced by students, educators, and software professionals who regularly work with hand-drawn flowcharts. By combining Google Gemini's multimodal capabilities with a purpose-built interactive editor, the project demonstrates that digitization can be reduced from 15–30 minutes to a matter of seconds, with high accuracy. Support for multiple input methods makes the system adaptable to diverse real-world workflows. Beyond individual productivity, the platform reduces the skill barrier associated with professional diagramming tools, making accurate flowchart generation accessible to users with minimal technical background. The open-source technology stack enables free adoption, extension, and future research in intelligent document processing.